\section{Results}\label{sec:results}

\sonicomment{For personal reference, the three key takeaways from the results tables:

1. 

}
In this section, we discuss the experimental results of \methodname. We begin by discussing overall performance on 3 benchmarks. Then, we discuss comparisons with existing approaches. 

\sonicomment{Shift to analysis section -- we don't know if this works yet}
Finally, we observe how \methodname affects downstream performance by aiding other agents as a subagent in solving SWE tasks.
\begin{table*}[t]
\caption{SWE-Bench Verified Results}
\centering
\renewcommand{\arraystretch}{1.3}
\resizebox{\textwidth}{!}{%
\begin{tabular}{c|c|c*{9}{|c}}
\toprule
\toprule
\multirow{2}{*}{\textbf{Scaffold}} & \multirow{2}{*}{\textbf{LLM}} & \multirow{2}{*}{\textbf{Training Algo.}} & \multicolumn{3}{|c}{\textbf{File-level}} & \multicolumn{3}{|c}{\textbf{Module-level}} & \multicolumn{3}{|c}{\textbf{Function-level}} \\
\cline{4-12}
& & & Recall (\%) & Precision (\%) & F1 (\%) & Recall (\%) & Precision (\%) & F1 (\%) & Recall (\%) & Precision (\%) & F1 (\%) \\
\midrule
\multicolumn{12}{c}{Closed-Source LLMs} \\
\midrule
RepoSearcher & Claude-3.7-Sonnet & - & 89.24 & 20.24 & 32.3 & - & - & - & 66.08 & 18.64 & 26.91 \\
\midrule
\multirow{3}{*}{RepoNavigator} & GPT-5-Chat & - & 58.17 & 61.87 & 58.88 & - & - & - & 30.42 & 34.56 & 31.17 \\
& Claude-3.7-Sonnet & - & 72.26 & 75.95 & 73.01 & - & - & - & 31.03 & 34.43 & 31.72 \\
& Claude-4.5-Sonnet & - & 80.68 & 81.92 & 79.94 & - & - & - & 43.97 & 45.76 & 43.62 \\
\midrule
\multicolumn{12}{c}{Open-Source LLMs} \\
\midrule
Agentless & \multirow{5}{*}{Qwen2.5-32B} & - & 78.93 & 25.6 & 35.38 & - & - & - & 40.97 & 24.07 & 27.33 \\
\cmidrule{1-1}\cmidrule{3-12}
CoSIL & & - & 83.5 & 19.34 & 30.77 & - & - & - & 55.38 & 14.85 & 22.11 \\
\cmidrule{1-1}\cmidrule{3-12}
LocAgent & & - & 79.39 & 34.18 & 44.18 & - & - & - & 46.79 & 16.29 & 21.48 \\
\cmidrule{1-1}\cmidrule{3-12}
OrcaLoca & & - & 59.57 & 59.51 & 58.11 & - & - & - & 39.14 & 25.59 & 28.72 \\
\midrule
\multirow{2}{*}{RepoSearcher} & Qwen2.5-7B & \multirow{2}{*}{Distillation+GRPO} & 83.11 & 18.8 & 30.09 & - & - & - & 62.38 & 17.68 & 25.57 \\
& Qwen2.5-32B &  & 88.59 & 20.24 & 32.25 & - & - & - & 68.55 & 19.36 & 28.03 \\
\midrule
\multirow{3}{*}{RepoNavigator} & Qwen2.5-7B & \multirow{3}{*}{GRPO} & 50.62 & 53.83 & 51.63 & - & - & - & 26.69 & 30.34 & 27.49 \\
& Qwen2.5-14B & & 61.6 & 58.97 & 58.9 & - & - & - & 31.02 & 30.08 & 29.23 \\
& Qwen2.5-32B & & 67.29 & 70.76 & 67.75 & - & - & - & 33.71 & 37.19 & 34.09 \\
\midrule
\multirow{5}{*}{OpenHands} & Qwen3-4B-Instruct & - & 53.34 & 49.69 & 49.73 & 20.15 & 19.86 & 19.32 & 13.74 & 14.17 & 13.27 \\
& Qwen3-14B & - & 71.2 & 36.49 & 43.13 & 33.04 & 20.4 & 22.86 & 23.58 & 14.51 & 16.08 \\
& Qwen3-32B & - & 70.41 & 58.63 & 60.88 & 40.27 & 35.57 & 35.95 & 28.69 & 25.94 & 25.77 \\
\cmidrule{2-12}
& \methodname-4B & \multirow{2}{*}{GSPO} & 67.74 & 71.53 & 68.52 & 44.97 & 49.7 & 45.97 & 35.72 & 40.71 & 36.78 \\
& \methodname-14B & & 68.69 & 71 & 68.57 & 50.88 & 53.71 & 50.88 & 40.27 & 43.74 & 40.32 \\
\bottomrule
\end{tabular}%
}
\label{tab:swe_bench_results}
\end{table*}
\subsection{\methodname as a highly competitive code-localization agent}

\methodname for both 4B and 14B sizes perform at a much higher F1 level on all three SWE-tasks. Notably, while lagging behind RepoNavigator-32B on file-level F1, \methodname achieves higher function-level F1

\sonicomment{We only lag for pro, not for verified and lite}

\subsection{}


\begin{table*}[t]
\caption{\methodname outperforms competitive baselines on SWE-Bench Pro for file, module, and function-level retrieval, with significantly higher F1 scores. We see particularly large gains in function-level localization highlighting the effectiveness of \methodname in locating relevant code snippets within files. All models from the Qwen2.5 family are instruct variants, and all the recall, precision, and F1 scores are in \%}
\centering
\renewcommand{\arraystretch}{1.6}
\resizebox{\textwidth}{!}{%
\begin{tabular}{c|c|c*{9}{|c}}
\toprule
\toprule
\multirow{2}{*}{\textbf{Scaffold}} & \multirow{2}{*}{\textbf{LLM}} & \multirow{2}{*}{\textbf{Training Algo.}} & \multicolumn{3}{|c}{\textbf{File-level}} & \multicolumn{3}{|c}{\textbf{Module-level}} & \multicolumn{3}{|c}{\textbf{Function-level}} \\
\cline{4-12}
& & & Rec. & Prec. & \textbf{F1} & Rec. & Prec. & \textbf{F1} & Rec. & Prec. & \textbf{F1} \\
\midrule
Agentless & \multirow{4}{*}{Qwen2.5-32B} & - & 43.07 & 13.89 & 20.07 & - & - & - & 11.08 & 7.31 & 7.98 \\
\cmidrule{1-1}\cmidrule{3-12}
CoSIL & & - & 48.87 & 14.03 & 20.95 & - & - & - & 14.03 & 6 & 7.67 \\
\cmidrule{1-1}\cmidrule{3-12}
LocAgent & & - & 25.73 & 0.38 & 19.77 & - & - & - & 8.72 & 0.17 & 4.3 \\
\cmidrule{1-1}\cmidrule{3-12}
RepoSearcher & & - & 9 & 2.52 & 3.81 & & & & 2.52 & 2.46 & 2.31 \\
\midrule
\multirow{3}{*}{RepoNavigator} & Qwen2.5-7B & \multirow{3}{*}{GRPO} & 36.36 & 48.13 & 39.74 & - & - & - & 12.33 & 21.26 & 14.29 \\
& Qwen2.5-14B & & 46.85 & 58.64 & 49.72 & - & - & - & 16.05 & 25.25 & 18.06 \\
& Qwen2.5-32B & & 53.49 & 68.69 & 57.57 & - & - & - & 18.13 & 29.44 & 20.72 \\
\midrule
\multirow{5}{*}{OpenHands} & Qwen3-4B-Instr. & - & 35.59 & 44.42 & 36.96 & 10.19 & 17.46 & 11.78 & 7.01 & 12.16 & 8.12 \\
& Qwen3-14B & - & 48.22 & 28.48 & 30.08 & 14.21 & 13.82 & 11.87 & 9.92 & 9.97 & 8.2 \\
& Qwen3-32B-think & - & 54.55 & 49.65 & 46.85 & 22.62 & 27.18 & 21.94 & 11.93 & 17.82 & 12.31 \\
\cmidrule{2-12}
& \methodname-4B & \multirow{2}{*}{GSPO} & 46.16 & 68.98 & 51.77 & 31.13 & 56.05 & 36.97 & 23.73 & 48.65 & 29.03 \\
& \methodname-14B & & 48.81 & 68.81 & 53.63 & 32.02 & 53.8 & 37.13 & 23.76 & 46.09 & 28.74 \\
\bottomrule
\end{tabular}%
}
\label{tab:swe_bench_pro_results}
\end{table*}


\autoref{tab:swe_bench_pro_results} presents the performance of various scaffolds and language models on the code localization task across three granularity levels. The baseline approaches; Agentless, CoSIL~\cite{}, LocAgent~\cite{}, and RepoSearcher~\cite{} using Qwen2.5-32B exhibit a pattern of high recall and poor precision at the file and function level. This suggests systematic over-retrieval, where models identify many candidate files but lack the discriminative capability to filter irrelevant results. 

% In contrast, 

The introduction of reinforcement learning-based training yields substantial improvements with approaches such as RepoNavigator~\cite{zhang2026toolenoughreinforcementlearning} with GRPO result in both higher precision and recall.

% \methodname models within the OpenHands framework, which achieve competitive results with significantly fewer parameters. CS-OSS-4B reaches 51.77 F1 at the file level and 29.03 F1 at the function level—the latter representing a 136% improvement over the untrained Qwen3-32B-think model (12.31 F1) despite being 8× smaller.
% The performance gap between trained and untrained models widens progressively at finer granularities. At the function level, where precise localization is most challenging, GSPO-trained models achieve F1 scores of 28.74–29.03 compared to 7.67–12.31 for untrained approaches. This disparity suggests that specialized training is essential for capturing fine-grained code structure and semantic relationships. Furthermore, the diminishing returns observed between CS-OSS-4B and CS-OSS-14B (29.03 vs. 28.74 function-level F1) indicate that GSPO training enables efficient utilization of model capacity, achieving near-optimal performance at smaller scales. 


\subsection{\methodname Improves Downstream Performance on SWE-tasks}

We further investigate how well \methodname aids in actual SWE tasks. To do this, we append CS-OSS prediction to the user message for another agent. We evaluate Qwen3-8B with on 100 samples of SWE-Bench Verified. Initially, Qwen3-8B has poor performance with only X\% resolve rate. However, when using \methodname, we observe a performance increase of X\%. 


